\documentclass[11pt]{article}
\usepackage[T1]{fontenc}
\usepackage[utf8]{inputenc}
\usepackage[a4paper,margin=1in]{geometry}
\usepackage{lmodern}
\usepackage{microtype}
\usepackage{amsmath,amssymb,mathtools}
\usepackage{enumitem}
\usepackage{tikz-cd}
\setlength{\parindent}{0pt}
\setlength{\parskip}{0.58em}
\allowdisplaybreaks
\newcommand{\T}{\mathcal T}
\newcommand{\A}{\mathcal A}
\newcommand{\Af}{\mathcal A_f}
\newcommand{\Rep}{\operatorname{Rep}}
\newcommand{\id}{\mathrm{id}}
\begin{document}

\begin{center}
{\large\bfseries Deligne's Letter to Vasiu}
\end{center}

\begin{flushright}
November 30, 2011
\end{flushright}

Dear Vasiu,

Let \(\T\) be a tannakian category over an algebraically closed field \(k\). Your question amounts to: does \(\T\) always have a fiber functor over \(k\)? I believe it is true. The following argument is presumably too pedestrian, hiding which ``compacity arguments'' are really relevant.

Let \(\A\) be the set of strictly full subcategories of \(\T\), stable by \(\otimes\), subquotients and duals. Let \(\Af\subset\A\) be the set of those generated by finitely many objects, using the same operations. For \(\alpha\in\A\), I denote by \(\T_\alpha\) the corresponding subcategory (as defined: \(\T_\alpha=\alpha\)). The set \(\A\) is ordered by inclusion. I assume we already know the existence and unicity up to isomorphism of fiber functors for the \(\T_\alpha\), \(\alpha\in\Af\).

If \(\T_\alpha\subset\T_\beta\), it makes sense to say that a fiber functor \(\omega_\beta\) on \(\T_\beta\) extends (on the nose) a fiber functor \(\omega_\alpha\) on \(\T_\alpha\). If an extension up to isomorphism exists, an actual extension exists too.

Let us order the set of pairs
\[
(\alpha,\omega_\alpha),\qquad \alpha\in\A,\quad \omega_\alpha\text{ a fiber functor on }\T_\alpha,
\]
by ``\(\omega_\beta\) extends \(\omega_\alpha\)''. To avoid set-theoretical difficulties, one could consider only the fiber functors \(\omega_\alpha\) such that the \(\omega_\alpha(X)\) take values in the set of vector spaces \(k^n\), \(n\in\mathbb N\).

The ordered set of \((\alpha,\omega_\alpha)\) is inductive (Bourbaki, \emph{Ens.}, Ch.~3, \S2, no.~4): if \(I\) is a totally ordered subset, the ``union'' of the \((\alpha,\omega_\alpha)\) is a \(\T_0\) in \(\A\),
\[
\T_0=\bigcup_{(\alpha,\omega_\alpha)\in I}\T_\alpha,
\]
with the fiber functor \(\omega_0\) characterized by \(\omega_0|_{\T_\alpha}=\omega_\alpha\). This ``union'' majorizes \(I\). By Bourbaki, loc. cit., Th.~2, the ordered set of the \((\alpha,\omega_\alpha)\) has a maximal element \((\T_1,\omega_1)\). To prove that \(\T_1=\T\), it suffices to prove the following.

\textbf{Lemma 1.} Let \(\T'\) be in \(\A\) and let \(\T''\) be in \(\Af\). Let \(\langle\T',\T''\rangle\) be the category generated by \(\T'\) and \(\T''\); it is in \(\A\). Then any fiber functor \(\omega'\) on \(\T'\) can be extended (on the nose, or up to isomorphism; this amounts to the same) to \(\langle\T',\T''\rangle\).

I first prove:

\textbf{Lemma 2.} Suppose \(\T'\) is in \(\Af\) too. ``Restriction'' is then an equivalence of categories
\[
\left\{
\begin{array}{c}
\text{fiber functors }\omega\text{ on }\langle\T',\T''\rangle
\end{array}
\right\}
\xrightarrow{\sim}
\left\{
\begin{array}{c}
\text{triples of a fiber functor }\omega'\text{ on }\T',\ \omega''\text{ on }\T'',\\
\text{and an isomorphism }\tau\text{ of the restrictions of }\omega'\text{ and }\omega''\\
\text{to }\T'\cap\T''\ \bigl(\T'\cap\T''\text{ is also in }\Af\bigr)
\end{array}
\right\}.
\]

We may assume that \(\langle\T',\T''\rangle\) is the category of representations of a group \(G\), and that, for invariant subgroups \(A\) and \(B\), \(\T'\) (resp. \(\T''\)) is the subcategory of representations where \(A\) (resp. \(B\)) acts trivially. That they generate \(\langle\T',\T''\rangle\) means that \(A\cap B=\{e\}\). The intersection \(\T'\cap\T''\) is the category of representations on which \(AB\) acts trivially.

The triples \((\omega',\omega'',\tau)\) are all isomorphic: as all \(\omega'\) (resp. all \(\omega''\)) are isomorphic, it suffices to see that \((\omega',\omega'',\tau_1)\) and \((\omega',\omega'',\tau_2)\) are isomorphic. Indeed, \(\tau_1\) and \(\tau_2\) differ by an automorphism of \(\omega'|_{\T'\cap\T''}\), and such an automorphism lifts to an automorphism of \(\omega'\):
\[
G/A(k)\longrightarrow G/AB(k)\quad\text{is onto.}
\]

We hence have here categories with just one isomorphism class of objects, and the question is to compare automorphism groups. We need to check
\[
G\xrightarrow{\sim}\left\{(g',g'')\in G/A\times G/B\mid g'\text{ and }g''\text{ have the same image in }G/AB\right\}.
\]
Equivalently, this is the Cartesian square
\[
\begin{tikzcd}[column sep=large,row sep=large]
G \arrow[r] \arrow[d] & G/A \arrow[d] \\
G/B \arrow[r] & G/AB .
\end{tikzcd}
\]
Injectivity of this morphism of groups amounts to \(A\cap B=\{e\}\). For surjectivity, if \(\widetilde g'\), \(\widetilde g''\) lift \(g'\), \(g''\), then, after writing \(\widetilde g''=\widetilde g'ab^{-1}\) with \(a\in A\), \(b\in B\), the element \(\widetilde g'a=\widetilde g''b\in G\) maps to \((g',g'')\).

\textbf{Lemma 3.} The same as Lemma 2, but with \(\T'\) only assumed to be in \(\A\).

Let \(\mathcal B\) be the set of the \(\T_\alpha\) (\(\alpha\in\Af\)) contained in \(\T'\). One has
\[
\langle\T',\T''\rangle=\bigcup_{\beta\in\mathcal B}\langle\T_\beta,\T''\rangle.
\]
One also has equivalences
\[
\left\{\text{fiber functors on }\T'\right\}
\xrightarrow{\sim}
\left\{
\begin{array}{c}
\text{fiber functors }\omega_\beta\text{ on the }\T_\beta,\text{ plus a compatible system}\\
\text{of isomorphisms }\omega_\beta|_{\T_\gamma}\xrightarrow{\sim}\omega_\gamma\text{ for }\T_\gamma\subset\T_\beta,\\
\text{with the compatibility condition for }\T_\delta\subset\T_\gamma\subset\T_\beta
\end{array}
\right\}.
\]
The same holds for fiber functors on \(\langle\T',\T''\rangle=\bigcup\langle\T_\beta,\T''\rangle\).

The \(\T_\beta\cap\T''\) have a largest element: if \(\T''=\Rep(G'')\), they correspond to invariant subgroups of \(G''\); subgroups are closed subschemes, and one uses the noetherian property. If \(\beta_0\) is such that \(\T_{\beta_0}\cap\T''\) is the largest \(\T_\beta\cap\T''\), then for any \(\beta\geq\beta_0\), extending \(\omega_\beta=\omega|_{\T_\beta}\) to \(\langle\T_\beta,\T''\rangle\) amounts to extending \(\omega_\beta|_{\T_\beta\cap\T''}\) to \(\T''\) (Lemma 2), that is, to extending \(\omega_{\beta_0}\) from \(\T_{\beta_0}\cap\T''\) to \(\T''\). If we choose one such extension, we get, up to unique isomorphism, a system of extensions of the \(\omega_\beta\) to \(\langle\T_\beta,\T''\rangle\), and by gluing them an extension of \(\omega\) to \(\langle\T',\T''\rangle\).

This is all we need to conclude that \(\T'=\T\).

The proof should be cleaned up. After all, we are proving that some projective system of gerbes (of the fiber functors on the \(\T_\alpha\), \(\alpha\in\Af\)), where ``projective system'' is taken in a 2-categorical sense, has a nonempty projective limit (again a limit in a 2-categorical sense). Maybe such a translation would make the ``compacity arguments'' used clearer. There were two of them: ``fiber functor is a property of finite type'' (cf. Bourbaki, \emph{Ens.}, Ch.~3, \S4, no.~5), and the noetherian property for subgroups.

The same arguments give unicity up to isomorphism of \(\omega\), over \(k\): we get a maximal \(\T_\alpha\) (\(\alpha\in\A\)) over which we have an isomorphism, and extend further if \(\T_\alpha\neq\T\).

Best,

P.~Deligne

\end{document}
